\chapter{Stage Appendix: Part VIII -- Relaxed Open-System Branch, Barrier Compiler, Export Kernel, and Material Thresholds}
\label{app:stage-part-8}

\section*{Stage range: 243--253}
\addcontentsline{toc}{section}{Stage range: 243--253}

This appendix expands the Part VIII theorem-block chapter into a stage-by-stage verification ledger.  The main-body chapter states the relaxed/open-system branch in citation-ready form; the present appendix preserves the derivational path from Stage~243 through Stage~253 and records the exact algebraic objects that a reader must check to audit the branch.

The result anchor served by this appendix is \resultanchor{MTDC-T12}.  Downstream documents should cite this anchor when they need the derivation of the relaxed-constraint branch declaration, the recovery map back to the strict front end, the selected leakage/work compiler, the non-rigid mouth and dressing packet, the compensated multimode source compiler, the relaxed stationary barrier law, the dynamic event-chain compiler, the helicity-export diagnostic, the crossing-versus-collapse window, the microscopic cubic/quintic export kernel, the vacuum/lattice heat partition, or the material-threshold companion.

The conceptual firewall is the most important part of this appendix.  The relaxed branch does \emph{not} undo the same-charge verdict from Part VII.  It does not produce a new long-range same-charge attraction, and it does not introduce a new linear dynamic kernel class.  The relaxed branch is a codimension-three lift of the strict selected-support/orbit packet.  Its standard-recovery slice is
\begin{equation}
\label{eq:app-part08-intro-recovery-slice}
\ell_w=0,
\qquad
f_U=0,
\qquad
a=b=0,
\end{equation}
which implies
\begin{equation}
\label{eq:app-part08-intro-recovery-observables}
S_{\rm leak}=0,
\qquad
\mathcal W_w=0,
\qquad
U=0,
\qquad
V=0,
\qquad
\mathcal D_{UV}=0,
\qquad
\varsigma\equiv 1.
\end{equation}
Away from this slice, the relaxed branch can reshape a near-contact barrier by leakage, scalar-photon work, non-rigid dressing drain, and compensated source response.  Those are branch-local open-system or support/dressing/source effects.  They do not alter the long-range kernel firewall.

The central stationary object of this part is the lowered front end
\begin{equation}
\label{eq:app-part08-intro-veff230}
V_{\rm eff}^{(247)}(r)
=
V_{\rm short}^{(1p)}(r)
-
\lambda_L S_{\rm leak}(r)
-
\lambda_W\mathcal W_w^{\rm sess}(r)
-
\Delta E_{UV}(r)
-
\mathcal M_{\sigma}(r).
\end{equation}
The central dynamic object is the event-chain energy law
\begin{equation}
\label{eq:app-part08-intro-event-energy}
E=\frac{m_s}{2}\dot r^{\,2}+V_{\rm eff}^{(247)}(r).
\end{equation}
The central microscopic export object is the super-Ohmic active-leg kernel
\begin{equation}
\label{eq:app-part08-intro-export-kernel}
K_{\rm exp}(s)=\Gamma_3s^3+\Gamma_5s^5.
\end{equation}
The central materials companion is the four-ratio screen
\begin{equation}
\label{eq:app-part08-intro-material-screen}
\Pi_{\rm ep}\ge 1,
\qquad
\Pi_{\chi}\ge 1,
\qquad
\Pi_k\ge 1,
\qquad
\Pi_T(T)\ge 1.
\end{equation}

\begin{longtable}{@{}L{0.075\textwidth}L{0.35\textwidth}L{0.20\textwidth}L{0.325\textwidth}@{}}
\caption{Part VIII appendix source and status map.}\label{tab:app-part08-source-status-map}\\
\toprule
Stage & Working title & Primary status & Main output \\
\midrule
\endfirsthead
\toprule
Stage & Working title & Primary status & Main output \\
\midrule
\endhead
243 & Relaxed-constraint branch declaration & Mixed: \StatusEffective{}, \StatusExactClosure{} & Codimension-three lift of the Stage 242 front end, exact recovery map \((\ell_w,f_U,a,b)=(0,0,0,0)\), leakage/work lane, non-rigid \((U,V)\) lane, compensated source lane, and strict short-range kernel-span theorem. \\
\addlinespace
244 & Selected leakage and scalar-photon work compiler & \StatusExactClosure{} & Exact Gaussian one-mode leakage law, exact bulk and session work scalars, selected-support pullback through \(\Pi_{\rm tr}\), support/orbit split, and orientation parity. \\
\addlinespace
245 & Non-rigid mouth and \(U/V\) drain packet & \StatusExactClosure{} & Exact finite compiler for \(\mathcal T^2\), \(\epsilon_\eta\), and \(R_{\rm target}\); non-rigid determinant \(\Delta_{UV}\); dependent microscopic correction; positive \(U/V\) drain; and first-order packet \((\Xi_1,\mathcal R_1)\). \\
\addlinespace
246 & Compensated multimode source compiler & \StatusExactClosure{} & Mean-preserving two-mode source family, exact sign-change test, mouth-bias and shell-loading moments, invertible two-moment source map, mixed-to-shell loading ratio, and transported Session-I readback. \\
\addlinespace
247 & Relaxed stationary barrier front end & Mixed: \StatusExactClosure{}, \StatusEffective{}, \StatusNumerical{} & One-port short-range baseline, relaxed lowering identity, stationary barrier compiler \(V_{\rm eff}^{(247)}\), and one-point Session-I benchmark decomposition. \\
\addlinespace
248 & Dynamic event-chain compiler & Mixed: \StatusExactClosure{}, \StatusReduced{}, \StatusNumerical{} & Energy integral, peak and threshold-speed compiler, turning points, WKB action, transmission ratio, dynamic diagnostics, and Session-II crossing/tunneling benchmark. \\
\addlinespace
249 & Helicity-export diagnostic & Mixed: \StatusExactClosure{}, \StatusEffective{}, \StatusNumerical{} & Exact projected subscale-helicity identity, aligned/anti-aligned orientation closure, peak and integrated ratio compilers, and Session-II hidden-channel export packet. \\
\addlinespace
250 & Crossing-vs-collapse Goldilocks window & Mixed: \StatusExactClosure{}, \StatusEffective{}, \StatusNumerical{} & Event-chain transit time, characteristic crossing time, dressing-leg collapse time, survival ratio, lower safe edge, speed-space compiler, heavy-throat scaling, and trigger-width sensitivity. \\
\addlinespace
251 & Microscopic damping/export kernel & Mixed: \StatusExactClosure{}, \StatusReduced{}, \StatusOpen{} & Derivative-coupled scalar cubic export, selected quadrupole quintic export, active-leg kernel \(K_{\rm exp}(s)=\Gamma_3s^3+\Gamma_5s^5\), exact event-safe half-plane, and channel-resolved power split. \\
\addlinespace
252 & Vacuum/lattice heat partition & Mixed: \StatusExactClosure{}, \StatusEffective{}, \StatusNumerical{}, \StatusOpen{} & Exact vacuum/lattice exported energies, one-scalar event-shape quotient, single-growth cold-event specialization, safe-edge exported-energy theorem, and event-equivalent export rates. \\
\addlinespace
253 & Physical calibration and material thresholds & Mixed: \StatusExactClosure{}, \StatusEffective{}, \StatusOpen{} & Physical calibration dictionary, coarse lattice-turnover projection factor, electron-phonon threshold, force-matched stiffness, Korringa ceiling, and four material-screening ratios. \\
\bottomrule
\end{longtable}

\claimstatus{Part VIII is a relaxed/open-system companion.  Its algebraic compilers are \StatusExactClosure{} inside the declared branch, one-mode, characteristic-timescale, single-growth, and physical-calibration closures.  The stationary and dynamic readbacks that use Session-I through Session-V benchmark numbers are \StatusNumerical{} insertions into exact formulas.  The branch declaration itself is an \StatusEffective{} extension around the strict Stage 242 front end, and any claim that the completed nonlinear moving-throat PDE actually realizes the relaxed packet, the microscopic export coefficients, or the physical material thresholds remains \StatusOpen{}.  No result in this part should be cited as a proof of a new same-charge long-range attraction or as a proof that the relaxed branch validates the strict conservative branch.}

\section{How to audit this appendix}
\label{sec:app-part08-audit-method}

A reader can verify the Part VIII branch by following four audit paths.
\begin{verificationchecklist}
  \item Check the recovery slice: impose \(\ell_w=f_U=a=b=0\) and verify that the lifted leakage, work, non-rigid, and source observables reduce to the strict Stage~242 values.
  \item Check the kernel firewall: verify that the one-port static and linear dynamic same-charge corrections stay in the short-range product span and never acquire a new \(1/r\) component.
  \item Check the compiler chain: insert the Stage~244 leakage/work packet, the Stage~245 \((U,V)\) packet, and the Stage~246 source packet into the Stage~247 stationary barrier, then pass that barrier through the Stage~248 event-chain and Stage~250 crossing/collapse tests.
  \item Check the physical companion: use Stage~251 to replace the phenomenological export envelope by \(K_{\rm exp}(s)=\Gamma_3s^3+\Gamma_5s^5\), use Stage~252 to split the exported energy into vacuum and lattice channels, and use Stage~253 to map the lattice channel into material-screening ratios.
\end{verificationchecklist}


\section{Block contract and theorem-level output}
\label{sec:app-part08-block-contract}

The contract of Part VIII is to keep three layers separate.  The first layer is the strict front end from Part VII: selected twin-support placement, rigid-mouth orbit lock, and the outgoing normalization packet.  The second layer is the relaxed branch around that front end: leakage, non-rigid mouth response, and compensated source shape.  The third layer is the physical companion: event-chain survival, export, heat partition, and material screening.  The chapter is useful only if those layers remain visibly separated.

The operational chain is
\begin{equation}
\label{eq:app-part08-operational-chain}
\begin{aligned}
\mathcal P_{242}
&\longrightarrow
\mathfrak B_{243}^{\rm relax}
\longrightarrow
(S_{\rm leak},\mathcal W_w,U,V,\mathcal M_\sigma)
\\
&\longrightarrow
V_{\rm eff}^{(247)}
\longrightarrow
(r_\pm,I_{\rm new},t_{\rm cross})
\\
&\longrightarrow
K_{\rm exp}(s)
\longrightarrow
(\Pi_{\rm ep},\Pi_\chi,\Pi_k,\Pi_T).
\end{aligned}
\end{equation}
The strict branch can be recovered from this chain by imposing \cref{eq:app-part08-intro-recovery-slice}.  The relaxed branch cannot be used to infer that the strict branch succeeds; it can only be used to analyze the open-system bypass companion.

\begin{theorem}[Part VIII relaxed-branch and export-companion theorem]
\label{thm:app-part08-main}
Inside the declared relaxed-constraint branch, the one-mode leakage/work closure, the non-rigid mouth/dressing closure, the compensated two-mode source closure, the reduced one-dimensional event-chain closure, the characteristic crossing/collapse closure, the microscopic cubic/quintic export closure, and the physical-calibration companion, the following statements hold.
\begin{enumerate}[leftmargin=2.2em]
  \item The relaxed branch is a codimension-three lift of the Stage 242 front end with exact recovery map
  \begin{equation}
  \label{eq:app-part08-main-recovery-map}
  \ell_w=0,
  \qquad
  f_U=0,
  \qquad
  a=b=0.
  \end{equation}
  On this slice all lifted leakage, work, non-rigid, and compensated-source observables reduce to the strict standard-recovery values.
  \item The same-charge static and linear dynamic kernel span remains short-ranged:
  \begin{equation}
  \label{eq:app-part08-main-short-range-span}
  \mathcal K_{\rm SR}
  =
  \operatorname{span}\left\{
  r^{-6},
  \frac{e^{-2\kappa r}}{r^4},
  \frac{e^{-4\kappa r}}{r^2}
  \right\}.
  \end{equation}
  In particular,
  \begin{equation}
  \label{eq:app-part08-main-short-range-limits}
  \lim_{r\to\infty}r\,\delta V_{\rm stat}(r)=0,
  \qquad
  \lim_{r\to\infty}r\,\Re\mathfrak V_{\rm dyn}(r,\omega)=0.
  \end{equation}
  \item On the selected-support branch, the leakage/work lane is explicit and support-side:
  \begin{equation}
  \label{eq:app-part08-main-leak-support}
  S_{\rm leak}(r)
  =
  \frac{8\sqrt2\,\eta_{\rm leak}\mu_w\rho_0}{\pi^{5/2}\lambda^3}
  \Lambda(r)\varrho(r),
  \end{equation}
  \begin{equation}
  \label{eq:app-part08-main-work-support}
  \mathcal W_w^{\rm sess}(r)
  =
  \frac{512\,\eta_{\rm leak}^{2}\mu_w q\rho_0}{\pi^4\lambda^2}
  \Lambda(r)^2\varrho(r)^2.
  \end{equation}
  These quantities are independent of the orbit-lock variables \((R_{\rm tr},R_{\rm target},\epsilon_\eta)\) within the declared pullback.
  \item The non-rigid mouth packet is exact in the physical logarithmic chart
  \begin{equation}
  \label{eq:app-part08-main-uv-chart}
  U=\ln\frac{\mathcal T^2}{\mathcal T_{\rm ref}^2},
  \qquad
  V=\ln\frac{\epsilon_\eta}{\epsilon_{\eta,{\rm ref}}}.
  \end{equation}
  With
  \begin{equation}
  \label{eq:app-part08-main-delta-uv}
  \Delta_{UV}=a_Ua_V-\chi_{UV}^{2},
  \end{equation}
  the stationary point is
  \begin{equation}
  \label{eq:app-part08-main-uv-stationary}
  U=\frac{a_Vf_U}{\Delta_{UV}},
  \qquad
  V=\frac{\chi_{UV}f_U}{\Delta_{UV}},
  \qquad
  \Delta_{UV}>0.
  \end{equation}
  The corresponding drain is nonnegative:
  \begin{equation}
  \label{eq:app-part08-main-uv-drain}
  \mathcal D_{UV}=\frac{\chi_{UV}^{2}a_Vf_U^2}{\Delta_{UV}^{2}}\ge 0.
  \end{equation}
  \item The first compensated two-mode source family is an exact two-moment compiler:
  \begin{equation}
  \label{eq:app-part08-main-source-family}
  \sigma_{a,b}(x)=1+a\cos(\pi x)+b\cos(2\pi x),
  \qquad
  \int_0^1\sigma_{a,b}(x)\,dx=1,
  \end{equation}
  with
  \begin{equation}
  \label{eq:app-part08-main-source-moments}
  \mathfrak g(a,b)=\frac{2}{\pi}\left(1+\frac{a}{3}-\frac{b}{15}\right),
  \qquad
  \mathcal S(a,b)=\frac{2\tanh(\pi/2)}{\pi}\left(1+\frac{a}{5}+\frac{b}{17}\right).
  \end{equation}
  Its normalized two-moment map has determinant \(14/425>0\), so the first compensated source family spans the two carried source moments exactly.
  \item The relaxed stationary barrier is exactly \cref{eq:app-part08-intro-veff230}, and the lowering identity is
  \begin{equation}
  \label{eq:app-part08-main-lowering-identity}
  V_{\rm short}^{(1p)}(r)-V_{\rm eff}^{(247)}(r)
  =
  \lambda_LS_{\rm leak}(r)
  +\lambda_W\mathcal W_w^{\rm sess}(r)
  +\Delta E_{UV}(r)
  +\mathcal M_\sigma(r).
  \end{equation}
  Hence, when the imported packet scalars and weights are nonnegative, \(V_{\rm eff}^{(247)}\le V_{\rm short}^{(1p)}\) pointwise.
  \item The dynamic event chain is determined by the conserved energy \cref{eq:app-part08-intro-event-energy}.  The lowered-branch classical threshold speed is
  \begin{equation}
  \label{eq:app-part08-main-vcrit-new}
  v_{\rm crit,new}=\sqrt{\frac{2}{m_s}\bigl(V_{\rm peak}-V(r_0)\bigr)},
  \end{equation}
  while the pure-Coulomb contact threshold is
  \begin{equation}
  \label{eq:app-part08-main-vcrit-coul}
  v_{\rm contact,Coul}
  =
  \sqrt{\frac{2}{m_s}\left(\frac1{r_{\rm contact}}-\frac1{r_0}\right)}.
  \end{equation}
  If \(v_{\rm crit,new}<v_0<v_{\rm contact,Coul}\), the lowered branch reaches contact while the pure Coulomb branch still turns back.
  \item The subbarrier WKB action is
  \begin{equation}
  \label{eq:app-part08-main-wkb-action}
  I_{\rm new}(E)=
  \frac{1}{\hbar_{\rm eff}}
  \int_{r_-(E)}^{r_+(E)}
  \sqrt{2m_s\bigl(V(r)-E\bigr)}\,dr,
  \end{equation}
  and the transmission ratio relative to Coulomb is
  \begin{equation}
  \label{eq:app-part08-main-transmission-ratio}
  \frac{T_{\rm new}(E)}{T_{\rm Coul}(E)}
  =
  \exp\bigl[-2\bigl(I_{\rm new}(E)-I_{\rm Coul}(E)\bigr)\bigr].
  \end{equation}
  \item The helicity-export comparison is diagnostic.  With orientation label \(\sigma=\pm1\), the reduced closure gives
  \begin{equation}
  \label{eq:app-part08-main-helicity-rate}
  \dot H_{{\rm sub},\sigma}=\Gamma_0(1+\sigma\alpha_h),
  \qquad
  |\alpha_h|<1
  \end{equation}
  for positive export in both branches.  The aligned branch dominates whenever \(\alpha_h>0\), but this does not by itself prove a microscopic spin sector.
  \item The crossing-versus-collapse safe edge is
  \begin{equation}
  \label{eq:app-part08-main-goldilocks-edge}
  E_{\rm edge}=V_{\rm peak}
  +
  \frac{m_s}{\mu_\eta}
  \frac{g_{UV}\chi_{\rm peak}\lambda_{\rm eff}^{2}}{2},
  \end{equation}
  equivalently
  \begin{equation}
  \label{eq:app-part08-main-goldilocks-speed}
  v_{\rm safe,min}^{2}
  =
  v_{\rm crit,new}^{2}
  +
  \frac{\lambda_{\rm eff}^{2}g_{UV}\chi_{\rm peak}}{\mu_\eta}.
  \end{equation}
  Under the declared characteristic closure this is a one-sided lower edge.
  \item The microscopic active-leg export kernel is super-Ohmic:
  \begin{equation}
  \label{eq:app-part08-main-kexp}
  K_{\rm exp}(s)=\Gamma_3s^3+\Gamma_5s^5,
  \qquad
  \Gamma_3,\Gamma_5\ge0.
  \end{equation}
  The positive growth root is reduced but not eliminated by any finite passive \((\Gamma_3,\Gamma_5)\) in this minimal closure.  The exact event-safe condition is
  \begin{equation}
  \label{eq:app-part08-main-safe-half-plane}
  \widehat\Gamma_3+s_c^2\widehat\Gamma_5
  \ge
  \frac{s_0^2-s_c^2}{s_c^3},
  \qquad
  \widehat\Gamma_n:=\Gamma_n/\mu_\eta.
  \end{equation}
  \item The heat partition is channel-resolved.  For event-shape quotient \(\mathfrak r_V=\mathcal I_3/\mathcal I_2\),
  \begin{equation}
  \label{eq:app-part08-main-heat-fraction}
  f_{\rm lat}(\mathfrak r_V)=
  \frac{\Gamma_3^{\rm lat}+\Gamma_5^{\rm lat}\mathfrak r_V}{\Gamma_3+\Gamma_5\mathfrak r_V},
  \end{equation}
  and on a single-growth cold event \(\mathfrak r_V=s^2\).  At the safe edge, the total minimum exported energy is
  \begin{equation}
  \label{eq:app-part08-main-safe-energy}
  E_{\rm exp,min}^{\rm safe}
  =
  \frac{V_{\rm in}^{2}}{2}(e^2-1)\mu_\eta(s_0^2-s_c^2).
  \end{equation}
  \item The physical-material companion is an exact calibration image of the Stage 252 event-equivalent rate.  With calibration factor \(\Upsilon_{\rm lat}>0\),
  \begin{equation}
  \label{eq:app-part08-main-ep-threshold}
  (\lambda_{\rm ep}\omega_D)_{\min}^{(253)}
  =
  \frac{f_{\rm lat}(s_c)\mu_\eta(s_0^2-s_c^2)}
  {\Upsilon_{\rm lat}\zeta_{\rm ep}s_ct_*}.
  \end{equation}
  Candidate hosts are screened by \cref{eq:app-part08-intro-material-screen}.
\end{enumerate}
Every item involving actual branch data, microscopic coefficients, physical units, or candidate material constants remains a realization or calibration question beyond the algebraic compiler.
\end{theorem}


\paragraph{Source layout.}
Stages~243--253 are now sourced from the canonical files \StageFile{stages/stage_243.tex} through \StageFile{stages/stage_253.tex}.  The raw Markdown notes and audit artifacts are tracked in a repo-local provenance index, so they do not have to appear in the reader-facing PDF.  Part VIII differs from Parts I--VII: these stage files carry the full derivation sections rather than compact stage cards.

\section{Stage 243: Relaxed constraint branch declaration}
\label{sec:app-part08:relaxed-constraint-branch-declaration}
\label{stage:243}
\stagefield{Verification}{SymPy audit: \StageFile{scripts/moving_throat_pde_stage243_relaxed_constraint_branch_declaration_and_short_range_open_system_compiler_sympy_audit.py}.  Mathematica audit: \StageFile{mathematica/moving_throat_pde_stage243_relaxed_constraint_branch_declaration_and_short_range_open_system_compiler_mathematica_audit.wl}.}


Stage 243 declares the relaxed branch without yet assembling a barrier.  Its purpose is to make the Session-I relaxations legal inside the derivation tree while preserving the same-charge short-range verdict.  The inherited base object is the Stage 242 front-end packet
\begin{equation}
\label{eq:app-part08-p225-packet}
\mathcal P_{242}
=
\Bigl(
\epsilon,
\varrho_{\rm phys},
\sigma_{\rm phys},
\text{ranking region},
R_{\rm tr},
R_{\rm target},
\epsilon_\eta,
 d\ln R_{\rm tr},
 d\ln R_{\rm target},
 d\ln\epsilon_\eta,
 N_Q-1
\Bigr),
\end{equation}
with selected-support packet
\begin{equation}
\label{eq:app-part08-s225-packet}
\mathcal S_{242}=(\zeta,M_{\rm mix},S(\zeta;\epsilon),M_{\rm tr}).
\end{equation}
The strict slice is
\begin{equation}
\label{eq:app-part08-std-branch-slice}
\mathfrak B_{\rm std}
=
\left\{
J^w=0,
\;U=0,
\;V=0,
\;\varsigma(z)\equiv 1
\right\}
\subset(\mathcal P_{242},\mathcal S_{242}).
\end{equation}
This definition is useful because it names exactly what is being relaxed: transverse-current suppression, rigid-mouth orbit lock, and positive-source Family-1 closure.

\subsection{Open-system leakage/work lane}
\label{subsec:app-part08-open-lane}

The open-system lift begins with the exact projected-continuity leakage identity
\begin{equation}
\label{eq:app-part08-stage243-leakage-identity}
S_{\rm leak}
=
-\bigl[W(w)j^w\bigr]_{-\infty}^{+\infty}
+
\int_{-\infty}^{+\infty}W'(w)j^w(w)\,dw.
\end{equation}
For the declaration-level Gaussian closure,
\begin{equation}
\label{eq:app-part08-stage243-gaussian-profile}
W(w)=\frac{e^{-w^2}}{\sqrt\pi},
\qquad
j^w(w)=\ell_w j_0\,w e^{-w^2},
\qquad
E_w(w)=E_0we^{-w^2},
\end{equation}
so the boundary term vanishes and
\begin{equation}
\label{eq:app-part08-stage243-leak-work}
S_{\rm leak}=-\frac{\sqrt2}{4}\ell_wj_0,
\qquad
\mathcal W_w=\frac{\sqrt{2\pi}}{8}\ell_wj_0E_0.
\end{equation}
The declaration packet for this lane is
\begin{equation}
\label{eq:app-part08-stage243-lw}
L_w=(S_{\rm leak},\mathcal W_w).
\end{equation}
Both components vanish on \(\ell_w=0\).

\subsection{Non-rigid mouth/dressing lane}
\label{subsec:app-part08-uv-lane}

The non-rigid lift uses the minimal free energy
\begin{equation}
\label{eq:app-part08-stage243-fuv}
\mathcal F_{UV}(U,V)
=
\frac12k_UU^2+\frac12k_VV^2-\chi_\lambda UV-f_UU.
\end{equation}
Stationarity gives
\begin{equation}
\label{eq:app-part08-stage243-uv-solution}
U=\frac{k_Vf_U}{k_Uk_V-\chi_\lambda^2},
\qquad
V=\frac{\chi_\lambda f_U}{k_Uk_V-\chi_\lambda^2},
\qquad
\frac{V}{U}=\frac{\chi_\lambda}{k_V}.
\end{equation}
The branch is admissible when
\begin{equation}
\label{eq:app-part08-stage243-uv-stability}
k_Uk_V-\chi_\lambda^2>0,
\end{equation}
and the drain is
\begin{equation}
\label{eq:app-part08-stage243-duv}
\mathcal D_{UV}=\chi_\lambda UV
=
\frac{\chi_\lambda^2k_Vf_U^2}{(k_Uk_V-\chi_\lambda^2)^2}\ge0.
\end{equation}
This lane is the finite relaxed continuation of the earlier weak-axisymmetric transfer-shape scalar: at infinitesimal order, \(\Xi_1=\delta U\).

\subsection{Compensated sign-changing source lane}
\label{subsec:app-part08-compensated-source-lane}

The compensated source lift is
\begin{equation}
\label{eq:app-part08-stage243-varsigma}
\varsigma(z)=1+a\cos(\pi z)+b\cos(2\pi z),
\qquad
z\in[0,1],
\end{equation}
with exact unit mean.  Setting \(y=\cos(\pi z)\),
\begin{equation}
\label{eq:app-part08-stage243-varsigma-quadratic}
\varsigma(y)=1-b+ay+2by^2,
\qquad
-1\le y\le 1.
\end{equation}
The interior stationary point and value are
\begin{equation}
\label{eq:app-part08-stage243-varsigma-min}
y_*=-\frac{a}{4b},
\qquad
\varsigma(y_*)=1-b-\frac{a^2}{8b},
\end{equation}
with the boundary candidates \(1+a+b\) and \(1-a+b\).  Thus the first source relaxation is not a vague sign-changing profile; it is an exact mean-preserving quadratic problem.

The full declaration packet is
\begin{equation}
\label{eq:app-part08-stage243-relaxed-branch}
\mathfrak B_{243}^{\rm relax}
=
\left\{(\mathcal P_{242},\mathcal S_{242});\;L_w,\;L_{UV},\;L_\varsigma\right\},
\end{equation}
where
\begin{equation}
\label{eq:app-part08-stage243-lanes}
L_w=(S_{\rm leak},\mathcal W_w),
\quad
L_{UV}=(U,V,\mathcal D_{UV}),
\quad
L_\varsigma=(\varsigma,\varsigma_{\min},\mathrm{signchg}).
\end{equation}

\subsection{Short-range theorem for the lifted branch}
\label{subsec:app-part08-short-range-theorem}

The most important Stage 243 output is the short-range firewall.  With primitive profiles
\begin{equation}
\label{eq:app-part08-stage243-source-profiles}
\mathcal S_Q(r)=r^{-3},
\qquad
\mathcal S_Y(r)=\frac{e^{-2\kappa r}}{r},
\end{equation}
the static one-port correction is
\begin{equation}
\label{eq:app-part08-stage243-static-correction}
\delta V_{\rm stat}(r)
=
-\frac12\left[
\frac{\mathcal C_6}{r^6}
+2\mathcal C_4\frac{e^{-2\kappa r}}{r^4}
+\mathcal C_2\frac{e^{-4\kappa r}}{r^2}
\right],
\end{equation}
and the linear monochromatic dynamic bundle has the same spatial span:
\begin{equation}
\label{eq:app-part08-stage243-dynamic-correction}
\Re\mathfrak V_{\rm dyn}(r,\omega)
=
-\frac12\left[
\frac{\mathcal C_6(\omega)}{r^6}
+2\mathcal C_4(\omega)\frac{e^{-2\kappa r}}{r^4}
+\mathcal C_2(\omega)\frac{e^{-4\kappa r}}{r^2}
\right].
\end{equation}
Thus the relaxed branch is a short-range/open-system bypass branch.  It can reshape near-contact energetics, but it cannot be cited as a new long-range same-charge force.


\section{Stage 244: Selected leakage and scalar-photon work compiler}
\label{sec:app-part08:selected-leakage-and-scalar-photon-work-compiler}
\label{stage:244}
\stagefield{Verification}{SymPy audit: \StageFile{scripts/moving_throat_pde_stage244_selected_branch_leakage_and_scalar_photon_work_compiler_sympy_audit.py}.  Mathematica audit: none yet.}


Stage 244 takes the declared \(J^w\neq0\) lane and attaches it to the selected-support packet.  The parent theory supplies both exact identities needed for the compiler: projected continuity gives \cref{eq:app-part08-stage243-leakage-identity}, and the localized Maxwell energy ledger gives
\begin{equation}
\label{eq:app-part08-stage244-em-energy}
\partial_tu_{\rm EM}+\partial_AS_{\rm EM}^{A}
=
-J^AE_A
=-(J^aE_a+J^wE_w).
\end{equation}
Thus \(S_{\rm leak}\) and \(J^wE_w\) are legitimate reduced observables once the mixed lane is opened.

\subsection{Gaussian one-mode leakage law}
\label{subsec:app-part08-stage244-gaussian-leakage}

The Stage 244 one-mode closure uses
\begin{equation}
\label{eq:app-part08-stage244-wlambda}
W_\lambda(w)=\frac{e^{-w^2/\lambda^2}}{\lambda\sqrt\pi},
\qquad
\phi_\lambda(w)=\frac{2w}{\sqrt\pi\lambda^3}e^{-w^2/\lambda^2},
\end{equation}
with
\begin{equation}
\label{eq:app-part08-stage244-ew-jw}
E_w(w;r)=-\mathcal E_0(r)\phi_\lambda(w),
\qquad
j^w(w;r)=\mu_w\rho_0E_w(w;r),
\qquad
J^w=qj^w.
\end{equation}
The exact projected leakage is
\begin{equation}
\label{eq:app-part08-stage244-sleak-e0}
S_{\rm leak}(r)
=
\frac{\sqrt2\mu_w\rho_0}{2\sqrt\pi\lambda^3}\mathcal E_0(r).
\end{equation}
The exact one-mode bulk work scalar is
\begin{equation}
\label{eq:app-part08-stage244-bulk-work}
\mathcal W_w^{\rm bulk}(r)
=
q\int j^wE_w\,dw
=
\frac{\sqrt2\mu_wq\rho_0}{2\sqrt\pi\lambda^3}\mathcal E_0(r)^2,
\end{equation}
so
\begin{equation}
\label{eq:app-part08-stage244-work-leak-relation}
\mathcal W_w^{\rm bulk}(r)=q\mathcal E_0(r)S_{\rm leak}(r).
\end{equation}
The Session-I scalar-photon work variable is the rescaled scalar
\begin{equation}
\label{eq:app-part08-stage244-session-work}
\mathcal W_w^{\rm sess}(r)
=
\frac{2\mu_wq\rho_0}{\lambda^2}\mathcal E_0(r)^2
=
2\sqrt{2\pi}\lambda\,\mathcal W_w^{\rm bulk}(r),
\end{equation}
or equivalently
\begin{equation}
\label{eq:app-part08-stage244-quadratic-law}
\mathcal W_w^{\rm sess}(r)
=
\frac{4\pi q\lambda^4}{\mu_w\rho_0}S_{\rm leak}(r)^2.
\end{equation}

\subsection{Selected-support pullback}
\label{subsec:app-part08-stage244-selected-pullback}

The selected-support demand from Stage 242 is
\begin{equation}
\label{eq:app-part08-stage244-pitr}
\Pi_{\rm tr}
=
\frac43C_{\rm mix}
=
\frac{32\Lambda(1-
\epsilon)}{3\pi^2}
=
\frac{16}{\pi^2}\Lambda\varrho,
\qquad
\varrho=\frac23(1-\epsilon).
\end{equation}
Stage 244 chooses the simplest amplitude pullback
\begin{equation}
\label{eq:app-part08-stage244-e0-pullback}
\mathcal E_0(r)=\eta_{\rm leak}\Pi_{\rm tr}(r).
\end{equation}
Substitution gives the selected-branch leakage law \cref{eq:app-part08-main-leak-support} and the selected-branch work law \cref{eq:app-part08-main-work-support}.  In the \((\Lambda,\epsilon)\) notation the same work law is
\begin{equation}
\label{eq:app-part08-stage244-work-epsilon}
\mathcal W_w^{\rm sess}(r)
=
\frac{2048\eta_{\rm leak}^{2}\mu_wq\rho_0}{9\pi^4\lambda^2}
\Lambda(r)^2(1-\epsilon(r))^2.
\end{equation}
The formulas depend only on selected-support variables.  Therefore,
\begin{equation}
\label{eq:app-part08-stage244-support-orbit-split}
\partial_{R_{\rm tr}}S_{\rm leak}
=
\partial_{R_{\rm target}}S_{\rm leak}
=
\partial_{\epsilon_\eta}S_{\rm leak}=0,
\end{equation}
and the same statement holds for both work scalars within the declared pullback.

The orientation parity is exact:
\begin{equation}
\label{eq:app-part08-stage244-parity}
S_{\rm leak}(-\eta_{\rm leak})=-S_{\rm leak}(\eta_{\rm leak}),
\qquad
\mathcal W_w(-\eta_{\rm leak})=\mathcal W_w(\eta_{\rm leak}).
\end{equation}
Thus the leakage sign is orientation-sensitive, but the exported-work magnitude is not.


\input{stages/stage_245}

\input{stages/stage_246}

\input{stages/stage_247}

\section{Stage 248: Dynamic event-chain compiler}
\label{sec:app-part08:dynamic-event-chain-compiler}
\label{stage:248}
\stagefield{Verification}{SymPy audit: \StageFile{scripts/moving_throat_pde_stage248_dynamic_event_chain_compiler_from_relaxed_stationary_barrier_front_end_turning_point_threshold_speed_and_wkb_sympy_audit.py}.  Mathematica audit: \StageFile{mathematica/moving_throat_pde_stage248_dynamic_event_chain_compiler_from_relaxed_stationary_barrier_front_end_turning_point_threshold_speed_and_wkb_mathematica_audit.wl}.}


Stage 248 promotes \(V_{\rm eff}^{(247)}\) into a reduced scattering/event-chain compiler.  Write
\begin{equation}
\label{eq:app-part08-stage248-vdef}
V(r):=V_{\rm eff}^{(247)}(r).
\end{equation}
The reduced radial equation is
\begin{equation}
\label{eq:app-part08-stage248-radial-eom}
m_s\ddot r=-V'(r),
\end{equation}
with conserved energy \cref{eq:app-part08-intro-event-energy}.  If the branch has a single barrier peak
\begin{equation}
\label{eq:app-part08-stage248-peak}
V'(r_{\rm peak})=0,
\qquad
V''(r_{\rm peak})<0,
\qquad
V_{\rm peak}=V(r_{\rm peak}),
\end{equation}
then the finite-radius threshold law is \cref{eq:app-part08-main-vcrit-new}.  The pure-Coulomb contact threshold is \cref{eq:app-part08-main-vcrit-coul}.  The classical comparison window is therefore
\begin{equation}
\label{eq:app-part08-stage248-classical-window}
v_{\rm crit,new}<v_0<v_{\rm contact,Coul}.
\end{equation}

For subbarrier energy \(V(r_0)<E<V_{\rm peak}\), turning points solve
\begin{equation}
\label{eq:app-part08-stage248-turning-points}
V(r_-(E))=E,
\qquad
V(r_+(E))=E,
\qquad
r_-(E)<r_+(E),
\end{equation}
and transport as
\begin{equation}
\label{eq:app-part08-stage248-turning-transport}
\frac{dr_\pm}{dE}=\frac1{V'(r_\pm(E))}.
\end{equation}
The WKB action and transmission ratio are \cref{eq:app-part08-main-wkb-action,eq:app-part08-main-transmission-ratio}.  The pure-Coulomb outer turning point is exact,
\begin{equation}
\label{eq:app-part08-stage248-coul-turn}
r_{\rm turn,Coul}(E)=\frac1E,
\end{equation}
and the Coulomb action closes as
\begin{equation}
\label{eq:app-part08-stage248-coul-action}
I_{\rm Coul}(E)
=
\frac{\sqrt{2m_s}}{\hbar_{\rm eff}}
\left[
\frac{\pi}{2\sqrt E}
-
\sqrt{r_{\rm contact}(1-Er_{\rm contact})}
-
\frac{\arcsin\sqrt{Er_{\rm contact}}}{\sqrt E}
\right]
\end{equation}
for \(Er_{\rm contact}<1\).

Near the top, with
\begin{equation}
\label{eq:app-part08-stage248-near-top}
V(r)=V_{\rm peak}-\frac{K_{\rm peak}}2(r-r_{\rm peak})^2+O((r-r_{\rm peak})^3),
\qquad
K_{\rm peak}=-V''(r_{\rm peak}),
\end{equation}
the leading top action is
\begin{equation}
\label{eq:app-part08-stage248-top-action}
I_{\rm top}(E)=
\frac{\pi(V_{\rm peak}-E)}{\hbar_{\rm eff}}
\sqrt{\frac{m_s}{K_{\rm peak}}}
+O((V_{\rm peak}-E)^{3/2}).
\end{equation}
The dynamic diagnostics carried forward are
\begin{equation}
\label{eq:app-part08-stage248-diagnostics}
\Xi_{\rm turn}(E)=\Xi_1(r_+(E)),
\qquad
\lambda_{\rm th}(E)=\left|\frac{E}{V'(r_+(E))}\right|.
\end{equation}
These tags are later used in the helicity and Goldilocks compilers; they do not alter the scalar event chain.


\input{stages/stage_249}

\input{stages/stage_250}

\section{Stage 251: Microscopic export kernel}
\label{sec:app-part08:microscopic-export-kernel}
\label{stage:251}
\stagefield{Verification}{SymPy audit: \StageFile{scripts/moving_throat_pde_stage251_microscopic_damping_export_kernel_replacing_the_phenomenological_v_leg_envelope_law_sympy_audit.py}.  Mathematica audit: none yet.}


Stage 251 replaces the phenomenological \(\gamma_{\rm tot}\dot V\) envelope law by the first microscopic odd export operator seen by the active dressing leg.  Project
\begin{equation}
\label{eq:app-part08-stage251-v-projection}
V(t)=\Pi_{V0}q_0(t)+\Pi_{V-}q_-(t)+\cdots,
\end{equation}
where \(q_0\) is the derivative-coupled scalar outlet and \(q_-\) is the selected mixed quadrupole outlet.

For the scalar lane, take
\begin{equation}
\label{eq:app-part08-stage251-scalar-coupling}
g_{W,0}(\omega)=\eta_0\omega,
\qquad
g_{A,0}(\omega)=0,
\qquad
\Pi_0^{\rm out}(\omega)=i\gamma_1\omega+O(\omega^3).
\end{equation}
With
\begin{equation}
\label{eq:app-part08-stage251-delta0}
\Delta_0=\Omega_{U,0}^2\Omega_{W,0}^2-R_0^2,
\end{equation}
the scalar outgoing factor gives
\begin{equation}
\label{eq:app-part08-stage251-gamma30}
\Gamma_{3,0}=\gamma_1\eta_0^2\frac{\Omega_{U,0}^4}{\Delta_0^2},
\qquad
\Gamma_3=\Pi_{V0}^{2}\Gamma_{3,0}.
\end{equation}
For the selected quadrupole outlet,
\begin{equation}
\label{eq:app-part08-stage251-gamma5}
\Gamma_{5,-}=\frac{a^5}{27c_s^5}P_{0,-},
\qquad
P_{0,-}=\frac{\beta_0s_-}{\lambda_-},
\qquad
\Gamma_5=\Pi_{V-}^{2}\Gamma_{5,-}.
\end{equation}
Thus the microscopic kernel is \cref{eq:app-part08-main-kexp}.

The active-leg equation is
\begin{equation}
\label{eq:app-part08-stage251-active-equation}
\mu_\eta\ddot V-\kappa_VV+
\Gamma_3V^{(3)}-
\Gamma_5V^{(5)}=
\mathcal S_V(t),
\qquad
\kappa_V=g_{UV}\chi_{\rm peak}.
\end{equation}
In Laplace form,
\begin{equation}
\label{eq:app-part08-stage251-laplace-d}
\mathcal D_V(s)=\mu_\eta s^2-\kappa_V+\Gamma_3s^3+\Gamma_5s^5.
\end{equation}
The passive power identity is
\begin{equation}
\label{eq:app-part08-stage251-power-identity}
\dot V(\Gamma_3V^{(3)}-\Gamma_5V^{(5)})
=
\frac{d}{dt}\mathcal S_{\rm odd}
-
\Gamma_3\ddot V^{2}
-
\Gamma_5\dddot V^{2},
\end{equation}
so the exported power is
\begin{equation}
\label{eq:app-part08-stage251-export-power}
\mathcal P_{\rm exp}=\Gamma_3\ddot V^{2}+\Gamma_5\dddot V^{2}\ge0.
\end{equation}

For homogeneous growth \(V=V_0e^{st}\), the characteristic polynomial is
\begin{equation}
\label{eq:app-part08-stage251-characteristic}
F(s)=\Gamma_5s^5+\Gamma_3s^3+\mu_\eta s^2-\kappa_V.
\end{equation}
For positive \((\Gamma_3,\Gamma_5,\mu_\eta,\kappa_V)\), \(F(0)<0\), \(F(s)\to+\infty\), and \(F'(s)>0\) for \(s>0\).  Therefore there is exactly one positive growth root.  The small-kernel slowdown is
\begin{equation}
\label{eq:app-part08-stage251-small-kernel}
s_+=s_0-\frac{\Gamma_3s_0^2+\Gamma_5s_0^4}{2\mu_\eta}+O(\Gamma^2),
\qquad
s_0=\sqrt{\kappa_V/\mu_\eta}.
\end{equation}
Thus finite passive cubic/quintic export slows but does not remove the unstable root in the minimal closure.

For an event with crossing rate \(s_c=1/t_{\rm cross}\), strict monotonicity of \(F\) gives the safe half-plane \cref{eq:app-part08-main-safe-half-plane}.  In the benchmark cold event, this becomes
\begin{equation}
\label{eq:app-part08-stage251-benchmark-half-plane}
\widehat\Gamma_3+0.3013336471\widehat\Gamma_5\ge 289.61004918.
\end{equation}
This is the exact event-level replacement for the old \(\gamma_{\rm safe}\) inequality.


\paragraph{Stage-pair note.}
Stages 252 and 253 form the physical companion to the export kernel. Stage 252 keeps the reduced microscopic language; Stage 253 maps it into physical threshold variables.

\input{stages/stage_252}

\input{stages/stage_253}

\section{Citation and downstream-use guardrails}
\label{sec:app-part08-downstream-guardrails}

Part VIII is useful because it lets downstream papers talk about the relaxed branch without contaminating the strict branch.  The safe citation language is:
\begin{quote}
The relaxed/open-system companion is a codimension-three lift of the strict selected-support/orbit front end.  It has an exact recovery map to the strict slice, preserves the no-new-long-range same-charge verdict, and supplies separate leakage/work, non-rigid dressing, compensated-source, barrier, dynamic-event, export-kernel, heat-partition, and material-screening compilers.
\end{quote}

The unsafe citation language is:
\begin{quote}
The relaxed barrier branch proves the strict selected same-charge branch, or supplies a new long-range attraction.
\end{quote}
That statement is false under the status discipline of this ledger.  The relaxed branch is a companion extension around the strict front end.  It is valuable precisely because it states what can change when open-system and non-rigid effects are allowed while preserving what cannot change: the same-charge long-range kernel verdict.

The active open tasks after this part are also explicit.  The completed moving-throat PDE still has to realize or reject the strict selected support/orbit/outgoing packet from Part VII.  If the relaxed companion is used, the actual branch must additionally return the leakage amplitude, non-rigid \((U,V)\) packet, compensated source packet, microscopic export coefficients \((\Gamma_3,\Gamma_5)\) and their vacuum/lattice split, the unit map \((t_*,\lambda_{\rm phys},E_*)\), and the coarse transport factor \(\Upsilon_{\rm lat}\).  Without those realized data, Part VIII remains an exact compiler rather than a completed physical-branch theorem.
